%! Factoring Trinomials — Unit 1, Lesson 1.4 — Group Activity
\documentclass[10pt]{article}
\usepackage{atda-article}
\usepackage{atda-boxes}
\newcommand{\TallMath}[1]{$\displaystyle #1\rule[-1.4em]{0pt}{3.2em}$}
\setlength{\workrowsep}{6pt}

\begin{document}

\pageheader{Unit 1, Lesson 1.4}{Group Activity}
% NAMESTRIP: no \namedateperiod here. The student writes their name once, on the
% cover this component is stapled behind. See references/conventions.md.

\vspace{0.15in}

\begin{scenariobox}[The Scenario]{royal}
\small
The art club is painting a \textbf{rectangular mural} on the cafeteria wall. The district approved
the job by \emph{area} only --- one polynomial, no drawing, no dimensions --- and the club has to
order border tape before it can start. Your job is to recover the shape from the polynomial.

Work with your group. \textbf{Everyone starts at Tier R.} When your whole group can do Tier R
without help, move on.
\end{scenariobox}

\vspace{0.10in}

\boxguard[16]
\begin{tcolorbox}[colback=white, colframe=black!40, title=\textbf{Tier R --- Remediate},
    fonttitle=\bfseries, arc=2mm, left=3mm, right=3mm, top=2mm, bottom=2mm, breakable]
\small
Five fluency moves. Every one of them shows up again in Tier A. Check each answer by multiplying
back.

\begin{enumerate}[leftmargin=1.6em, itemsep=4pt, topsep=4pt]
  \item Factor: $x^2+7x+12$
\begin{work}
  3 \cdot 4 &= 12, \quad 3+4 = 7 \\
  x^2+7x+12 &= (x+3)(x+4)
\end{work}

  \item Factor: $x^2-9x+20$
\begin{work}
  (-4)(-5) &= 20, \quad -4+(-5) = -9 \\
  x^2-9x+20 &= (x-4)(x-5)
\end{work}

  \item Factor: $x^2+2x-24$
\begin{work}
  6(-4) &= -24, \quad 6+(-4) = 2 \\
  x^2+2x-24 &= (x+6)(x-4)
\end{work}

  \item Factor: $x^2-3x-28$
\begin{work}
  (-7)(4) &= -28, \quad -7+4 = -3 \\
  x^2-3x-28 &= (x-7)(x+4)
\end{work}

  \item Factor: $2x^2+9x+4$ \quad (leading coefficient is not $1$ --- use $ac$)
\begin{work}
  ac &= 2 \cdot 4 = 8, \quad 1+8 = 9 \\
  2x^2+9x+4 &= 2x^2+x+8x+4 \\
            &= x(2x+1) + 4(2x+1) \\
            &= (2x+1)(x+4)
\end{work}
\end{enumerate}
\end{tcolorbox}

\vspace{0.10in}

\boxguard[16]
\begin{tcolorbox}[colback=white, colframe=black!40,
    title=\textbf{Tier A --- Approaching Proficiency},
    fonttitle=\bfseries, arc=2mm, left=3mm, right=3mm, top=2mm, bottom=2mm, breakable]
\small
The approved mural has area
\[ A(x) = 6x^2+17x+5 \quad \text{square feet.} \]

\begin{enumerate}[leftmargin=1.6em, itemsep=4pt, topsep=4pt]
  \item Factor $A(x)$ by splitting the middle term.
\begin{work}
  ac &= 6 \cdot 5 = 30, \quad 2+15 = 17 \\
  6x^2+17x+5 &= 6x^2+2x+15x+5 \\
             &= 2x(3x+1) + 5(3x+1) \\
             &= (3x+1)(2x+5)
\end{work}

  \item The club is told the mural stands $(3x+1)$ feet tall. Which factor is its width, and how do
        you know?

\writeline

  \item The wall allows $x=3$. Find the height, the width, the area, and the length of border tape
        the club must order --- then check your area against the original polynomial.
\begin{work}
  3x+1 &= 10 \text{ ft} \qquad 2x+5 = 11 \text{ ft} \\
  10 \cdot 11 &= 110 \text{ sq ft} \\
  6(9)+17(3)+5 &= 110 \text{ sq ft} \quad\checkmark \\
  \text{tape} &= 2(10)+2(11) = 42 \text{ ft}
\end{work}

  \item A second wall is offered with area $A(x)=10x^2+29x+21$. Factor it and give both
        dimensions.
\begin{work}
  ac &= 10 \cdot 21 = 210, \quad 14+15 = 29 \\
  10x^2+29x+21 &= 10x^2+14x+15x+21 \\
               &= 2x(5x+7) + 3(5x+7) \\
               &= (5x+7)(2x+3)
\end{work}

  \item A third wall is offered as $A(x)=4x^2+22x+24$. Factor it \textbf{completely} --- look at
        what all three terms share before you compute $ac$.
\begin{work}
  4x^2+22x+24 &= 2\left(2x^2+11x+12\right) \\
              &= 2\left(2x^2+3x+8x+12\right) \\
              &= 2\left[x(2x+3) + 4(2x+3)\right] \\
              &= 2(2x+3)(x+4)
\end{work}
\end{enumerate}
\end{tcolorbox}

\vspace{0.10in}

\boxguard[30]
\begin{tcolorbox}[colback=white, colframe=black!40, title=\textbf{Tier E --- Extension},
    fonttitle=\bfseries, arc=2mm, left=3mm, right=3mm, top=2mm, bottom=2mm, breakable]
\small

\begin{enumerate}[leftmargin=1.6em, itemsep=4pt, topsep=4pt]
  \item \textbf{Critique the work.} A student turned in these two answers. Neither is acceptable.
        Write the correct factorization for each, then name the mistake.
        \[ 3x^2+15x+18 = (3x+6)(x+3) \qquad\text{and}\qquad x^2-x-12 = (x-4)(x-3) \]
\begin{work}
  3x^2+15x+18 &= 3\left(x^2+5x+6\right) = 3(x+2)(x+3) \\
  x^2-x-12 &= (x-4)(x+3)
\end{work}
\writelines{2}

  \item \textbf{A negative out front.} A design is quoted as $A(x) = -2x^2+2x+24$ because the
        planner measured the wall from the far edge. Factor it completely by pulling out the
        negative GCF first.
\begin{work}
  -2x^2+2x+24 &= -2\left(x^2-x-12\right) \\
              &= -2(x-4)(x+3)
\end{work}

  \item \textbf{Justify.} A supplier claims every quadratic trinomial splits into two whole-number
        panels. The club asks for a panel of area $x^2+4x+6$ square feet with sides $(x+p)$ and
        $(x+q)$ for integers $p$ and $q$. Show that this order cannot be filled, and explain why
        checking a handful of pairs is enough to be certain.
\begin{work}
  1 \cdot 6 &= 6, \quad 1+6 = 7 \\
  2 \cdot 3 &= 6, \quad 2+3 = 5 \\
  (-1)(-6) &= 6, \quad -1+(-6) = -7 \\
  (-2)(-3) &= 6, \quad -2+(-3) = -5
\end{work}
\writelines{2}
\end{enumerate}
\end{tcolorbox}

\end{document}
